\chapter{INTRODUCTION}\label{chap-intro}
\thispagestyle{empty}

\hspace{1cm} Plasma is widely known as the fourth state of matter. It is formed when a gas receives enough energy for its atoms to become ionized. During this process, electrons are separated from atoms, creating a mixture of free electrons, ions, and neutral particles. Because of this composition, plasma behaves differently from solids, liquids, and gases. It has the ability to conduct electricity, respond to electric and magnetic fields, and produce chemically reactive species. These unique characteristics make plasma valuable in many scientific and industrial applications.

One important type of plasma is cold plasma, also called non-thermal plasma. In this form, the electrons have high energy, but the overall temperature of the gas remains relatively low. This property allows cold plasma to interact with surfaces and materials without causing significant heat damage. Cold plasma is typically generated at atmospheric pressure by applying strong electric fields to gases such as argon, helium, or even air. Due to its low temperature and reactive nature, cold plasma has gained attention for various modern applications.

This project focuses on the design and development of a Cold Plasma Generator, aiming to harness the unique properties of cold plasma for potential applications in fields such as medicine, environmental treatment, agriculture, and material processing. Unlike thermal plasma, cold plasma allows for the generation of reactive species without significantly raising the temperature of the target material, making it both efficient and versatile.

Through this project, we aim to understand the fundamental principles of plasma generation, explore the mechanisms behind cold plasma behavior, and demonstrate a working prototype of a cold plasma generator. The final goal is to evaluate its effectiveness and explore its real-world applications, particularly in sterilization, surface modification, and pollutant degradation.  

\section{Problem Definition}
The generation of stable cold plasma at atmospheric pressure requires the design of a compact, high-voltage power supply capable of delivering tens of kilovolts. Conventional plasma generators are often bulky, expensive, or unsuitable for laboratory-scale or portable applications. There is a need for a reliable system that can provide:  
\begin{itemize}
    \item A stable output of around 30~kV DC from a standard 230~V AC supply.  
    \item Voltage and power control to enable tunable plasma generation.  
    \item Integration with a mechanical setup to sustain cold plasma jets using carrier gases such as argon or air.  
\end{itemize}

\section{Objectives of the Project Work}
In this section, we outline the major objectives of the project. These objectives serve as the guiding framework for the design and development of the cold plasma generator.

The main objectives of this project are as follows:

\begin{enumerate}
    \item To design and implement a high-voltage plasma generator suitable for producing cold plasma at atmospheric pressure.
    \item To generate a stable output of approximately 30~kV DC using a high-voltage pulsed power supply and Cockcroft–Walton voltage multiplier.
    \item To achieve controllable output voltage levels by varying the input pulse width modulation (PWM) parameters through the STM32G431-Long Core Demo Board microcontroller from ST Electronics.
    \item To integrate the developed high-voltage power supply with the mechanical setup in order to produce and sustain a visible cold plasma jet using argon gas.
    \item To evaluate the potential applications of the generated cold plasma in areas such as sterilization, pollutant degradation, and surface modification.
\end{enumerate}

\section{Scope of the Project Work}
The scope of this project is limited to the design, fabrication, and testing of a high-voltage plasma generator at laboratory scale. The outcomes will be constrained by the following boundaries:
\begin{enumerate}
    \item The system will be designed to convert 230~V AC mains supply into approximately 30~kV DC at a switching frequency of 100~kHz, with an output power range of 150--250~W.

    \item The usability of the prototype is limited to laboratory testing and small-scale applications in medical sterilization, surface treatment, and environmental research.
\end{enumerate}
